\documentclass{article}
\usepackage{amsmath, amssymb}
\usepackage[utf8]{inputenc}
\usepackage[T1]{fontenc}

\begin{document}

\section*{Calcul de l'intégrale $I$}

Nous cherchons à évaluer l'intégrale suivante :
\[
I = \int_0^1 \frac{\ln(t) \ln^2(1-t)}{t} \, dt
\]

Considérons le changement de variable $x = 1 - t$. Ainsi, $dt = -dx$, les bornes d'intégration s'inversent de $1$ à $0$, et nous obtenons :
\[
I = \int_0^1 \frac{\ln(1-x) \ln^2(x)}{1-x} \, dx
\]

Nous savons que le développement en série entière de $-\frac{\ln(1-x)}{1-x}$ est généré par les nombres harmoniques $H_n = \sum_{k=1}^n \frac{1}{k}$ :
\[
- \frac{\ln(1-x)}{1-x} = \sum_{n=1}^\infty H_n x^n
\]

En appliquant ce développement à notre intégrale, nous avons :
\[
I = \int_0^1 \ln^2(x) \left( - \sum_{n=1}^\infty H_n x^n \right) \, dx
\]

Par la convergence uniforme (ou le théorème de convergence dominée de Lebesgue) de la série, nous pouvons intervertir l'intégrale et la somme :
\[
I = - \sum_{n=1}^\infty H_n \int_0^1 x^n \ln^2(x) \, dx
\]

Calculons maintenant l'intégrale intérieure $\int_0^1 x^n \ln^2(x) \, dx$. Nous savons, par exemple en utilisant la fonction Gamma ou des intégrations par parties successives, que :
\[
\int_0^1 x^n \ln^k(x) \, dx = (-1)^k \frac{k!}{(n+1)^{k+1}}
\]
Pour $k=2$, cela donne :
\[
\int_0^1 x^n \ln^2(x) \, dx = \frac{2}{(n+1)^3}
\]

En insérant ce résultat dans notre expression de $I$ :
\[
I = - \sum_{n=1}^\infty H_n \frac{2}{(n+1)^3} = -2 \sum_{n=1}^\infty \frac{H_n}{(n+1)^3}
\]

En effectuant le décalage d'indice $k = n + 1$, la série devient :
\[
I = -2 \sum_{k=2}^\infty \frac{H_{k-1}}{k^3}
\]
Puisque $H_0 = 0$, la sommation peut commencer à $k=1$ :
\[
I = -2 \sum_{k=1}^\infty \frac{H_{k-1}}{k^3}
\]

En exploitant la relation de récurrence des nombres harmoniques $H_{k-1} = H_k - \frac{1}{k}$, nous pouvons séparer cette somme :
\[
I = -2 \sum_{k=1}^\infty \frac{H_k - \frac{1}{k}}{k^3} = -2 \sum_{k=1}^\infty \frac{H_k}{k^3} + 2 \sum_{k=1}^\infty \frac{1}{k^4}
\]

Nous devons évaluer deux séries :
1. D'après Euler, la somme eulérienne de poids 4, $\sum_{k=1}^\infty \frac{H_k}{k^3}$, a pour valeur $\frac{5}{4} \zeta(4)$.
2. La seconde série correspond à la fonction zêta de Riemann, $\sum_{k=1}^\infty \frac{1}{k^4} = \zeta(4)$.

En substituant ces valeurs dans l'expression de l'intégrale, on obtient :
\[
I = -2 \left( \frac{5}{4} \zeta(4) \right) + 2 \zeta(4) = -\frac{5}{2} \zeta(4) + 2 \zeta(4) = - \frac{1}{2} \zeta(4)
\]

Sachant que $\zeta(4) = \frac{\pi^4}{90}$, nous concluons finalement :
\[
I = - \frac{1}{2} \left( \frac{\pi^4}{90} \right) = - \frac{\pi^4}{180}
\]

\vspace{0.5cm}
\subsection*{Méthode alternative : Produit de Cauchy (approache Math-OS)}

Une autre méthode, proposée par l'article de Math-OS, consiste à conserver l'intégrale sous sa forme initiale et à utiliser le produit de Cauchy pour le terme $\ln^2(1-t)$.

Considérons le développement en série entière du logarithme :
\[
-\ln(1-t) = \sum_{n=1}^\infty \frac{t^n}{n}
\]
En élevant cette série au carré par produit de Cauchy, on obtient :
\[
\ln^2(1-t) = \left( \sum_{i=1}^\infty \frac{t^i}{i} \right) \left( \sum_{j=1}^\infty \frac{t^j}{j} \right) = \sum_{n=2}^\infty \left( \sum_{k=1}^{n-1} \frac{1}{k(n-k)} \right) t^n
\]
On utilise la décomposition en éléments simples : $\frac{1}{k(n-k)} = \frac{1}{n} \left( \frac{1}{k} + \frac{1}{n-k} \right)$. 
La somme interne devient alors :
\[
\sum_{k=1}^{n-1} \frac{1}{k(n-k)} = \frac{1}{n} \sum_{k=1}^{n-1} \left( \frac{1}{k} + \frac{1}{n-k} \right) = \frac{2}{n} \sum_{k=1}^{n-1} \frac{1}{k} = \frac{2 H_{n-1}}{n}
\]
Ainsi, le développement en série s'écrit :
\[
\ln^2(1-t) = 2 \sum_{n=2}^\infty \frac{H_{n-1}}{n} t^n
\]

En substituant ce développement dans l'intégrale initiale $I$ :
\[
I = \int_0^1 \frac{\ln(t)}{t} \left( 2 \sum_{n=2}^\infty \frac{H_{n-1}}{n} t^n \right) \, dt = 2 \sum_{n=2}^\infty \frac{H_{n-1}}{n} \int_0^1 t^{n-1} \ln(t) \, dt
\]
Le calcul de l'intégrale $\int_0^1 t^{n-1} \ln(t) \, dt$ par intégration par parties donne :
\[
\int_0^1 t^{n-1} \ln(t) \, dt = \left[ \frac{t^n}{n} \ln(t) \right]_0^1 - \int_0^1 \frac{t^{n-1}}{n} \, dt = - \frac{1}{n^2}
\]

En réinjectant ce résultat, on retombe exactement sur la série exprimée avec les nombres harmoniques :
\[
I = 2 \sum_{n=2}^\infty \frac{H_{n-1}}{n} \left( - \frac{1}{n^2} \right) = -2 \sum_{n=2}^\infty \frac{H_{n-1}}{n^3} = -2 \sum_{k=1}^\infty \frac{H_k}{(k+1)^3}
\]
Ce résultat confirme que nous aboutissons à la même évaluation des séries eulériennes que dans la première méthode, menant au même résultat final $I = -\frac{\pi^4}{180}$.

\vspace{1cm}

\section*{Calcul de l'intégrale paramétrée $A_{p,q}$}

Soit $(p, q) \in \mathbb{N} \times \mathbb{N}^*$. L'intégrale que nous cherchons à évaluer est :
\[
A_{p,q} = \int_0^1 \frac{\ln^p(t) \ln(1-t^q)}{t} \, dt
\]

\subsection*{Étude de la convergence}

L'intégrande $f(t) = \frac{\ln^p(t) \ln(1-t^q)}{t}$ est défini et continu sur l'intervalle $]0, 1[$. Il y a deux problèmes potentiels aux bornes $0$ et $1$.

\textbf{En $t \to 0^+$ :}
On a l'équivalent $\ln(1-t^q) \sim -t^q$. Ainsi, au voisinage de $0$, le comportement asymptotique de l'intégrande est :
\[
f(t) \sim -\frac{\ln^p(t) t^q}{t} = -t^{q-1} \ln^p(t)
\]
Puisque $q \ge 1$, nous avons $q-1 \ge 0$. Par croissances comparées, la limite de $t^{q-1} \ln^p(t)$ en $0$ dépend de $q$ :
- Si $q > 1$, $\lim_{t \to 0} t^{q-1} \ln^p(t) = 0$. La fonction se prolonge par continuité en $0$, donc l'intégrale est faussement impropre (elle converge) en $0$.
- Si $q = 1$, l'équivalent est $-\frac{\ln^p(t)}{t} \times t = -\ln^p(t)$, qui est une fonction connue pour être intégrable sur $]0, a]$ (par exemple, car $t^{1/2} \ln^p(t) \to 0$).
Dans tous les cas, l'intégrale converge en $0$.

\textbf{En $t \to 1^-$ :}
Posons $t = 1 - h$ avec $h \to 0^+$. Alors $\ln(t) = \ln(1-h) \sim -h$.
De plus, $1 - t^q = 1 - (1-h)^q \sim qh$. 
Ainsi, l'intégrande se comporte comme :
\[
f(1-h) = \frac{\ln^p(1-h) \ln(1-(1-h)^q)}{1-h} \sim \frac{(-h)^p \ln(qh)}{1} = (-1)^p h^p \ln(qh)
\]
Puisque $p \in \mathbb{N}$, $p \ge 0$. Par croissances comparées, $\lim_{h \to 0} h^p \ln(qh) = 0$.
L'intégrande se prolonge donc par continuité en $1$ (la limite y est nulle). L'intégrale est faussement impropre en $1$.

En conclusion, l'intégrale $A_{p,q}$ est convergente pour tout $(p, q) \in \mathbb{N} \times \mathbb{N}^*$.

\vspace{0.5cm}
\subsection*{Évaluation de l'intégrale}

Nous cherchons maintenant à évaluer cette expression. Effectuons tout d'abord le changement de variable $u = t^q$. On a alors $t = u^{1/q}$, d'où $\ln(t) = \frac{1}{q} \ln(u)$, et $dt = \frac{1}{q} u^{\frac{1}{q}-1} \, du$.
L'élément différentiel devient :
\[
\frac{dt}{t} = \frac{\frac{1}{q} u^{\frac{1}{q}-1}}{u^{1/q}} \, du = \frac{du}{q u}
\]
Les bornes d'intégration restent inchangées ($t=0 \implies u=0$ et $t=1 \implies u=1$). 
En remplaçant dans l'intégrale, on obtient :
\[
A_{p,q} = \int_0^1 \frac{\left( \frac{1}{q} \ln(u) \right)^p \ln(1-u)}{q u} \, du = \frac{1}{q^{p+1}} \int_0^1 \frac{\ln^p(u) \ln(1-u)}{u} \, du
\]

Nous utilisons maintenant le développement en série entière de la fonction logarithme :
\[
\ln(1-u) = - \sum_{n=1}^\infty \frac{u^n}{n} \quad \text{pour } |u| < 1
\]
En substituant cette série dans notre intégrale, nous avons :
\[
A_{p,q} = \frac{1}{q^{p+1}} \int_0^1 \frac{\ln^p(u)}{u} \left( - \sum_{n=1}^\infty \frac{u^n}{n} \right) \, du
\]

Par application du théorème de convergence dominée (ou convergence monotone pour les fonctions de signe constant, selon la parité de $p$), il est justifié d'intervertir la somme et l'intégrale :
\[
A_{p,q} = - \frac{1}{q^{p+1}} \sum_{n=1}^\infty \frac{1}{n} \int_0^1 u^{n-1} \ln^p(u) \, du
\]

Évaluons l'intégrale $I_{n,p} = \int_0^1 u^{n-1} \ln^p(u) \, du$.
Posons le changement de variable $x = -\ln(u)$, avec $u = e^{-x}$ et $du = -e^{-x} \, dx$.
L'intervalle d'intégration devient $[+\infty, 0]$, ce qui donne :
\[
I_{n,p} = \int_{+\infty}^0 (e^{-x})^{n-1} (-x)^p (-e^{-x}) \, dx = \int_0^\infty e^{-nx} (-1)^p x^p \, dx = (-1)^p \int_0^\infty x^p e^{-nx} \, dx
\]
Avec un second changement de variable $y = nx$, on obtient $dx = \frac{dy}{n}$ et l'intégrale s'écrit :
\[
I_{n,p} = (-1)^p \int_0^\infty \left(\frac{y}{n}\right)^p e^{-y} \frac{dy}{n} = \frac{(-1)^p}{n^{p+1}} \int_0^\infty y^p e^{-y} \, dy
\]
On reconnaît l'intégrale eulérienne définissant la fonction Gamma : $\int_0^\infty y^p e^{-y} \, dy = \Gamma(p+1) = p!$.
Ainsi,
\[
I_{n,p} = \frac{(-1)^p p!}{n^{p+1}}
\]

En réinsérant ce résultat dans la somme pour $A_{p,q}$, on a :
\[
A_{p,q} = - \frac{1}{q^{p+1}} \sum_{n=1}^\infty \frac{1}{n} \left( \frac{(-1)^p p!}{n^{p+1}} \right) = \frac{(-1)^{p+1} p!}{q^{p+1}} \sum_{n=1}^\infty \frac{1}{n^{p+2}}
\]

Nous reconnaissons dans cette série la fonction zêta de Riemann, évaluée en $s = p+2$, d'expression $\zeta(s) = \sum_{n=1}^\infty \frac{1}{n^s}$.
Nous aboutissons donc à l'expression finale de l'intégrale paramétrée :
\[
A_{p,q} = \frac{(-1)^{p+1} p!}{q^{p+1}} \zeta(p+2)
\]

\end{document}
